% !PRINTED 422
\weberpage{422}

%<<CONT
coefficients of $x_1^{11}x_2,\ x_1^{22}x_2^2,\ x_1^{30}$ on both sides of the equations (14), and thus finds the relation
\weq{15}{t^5 - 10ft^3 + 45f^2t - T = 0,}
or also
\weq{16}{t^2(t^4 - 10ft^2 + 45f^2)^2 - T^2 = 0.}

From (16) and (15) there now result, by (10) and (12), the equations of the fifth degree for $r$ and for $u$
\weq{17}{r(r^2 - 10r + 45)^2 = z_1,}
\weq{18}{u^5 - \frac{10u^3}{z_1} + \frac{45u}{z_1^2} - \frac{1}{z_1^2} = 0.}

The last of these equations passes, by the substitution
\weq{19}{y = -\frac{1}{3u},\quad z_1 = 27\gamma}
into the form
\weq{20}{y^5 + 15y^4 - 10\gamma y^2 + 3\gamma^2 = 0}
which we have already come to know in \S.\,75 of the first volume as a normal form of the equation of the fifth degree.

\begin{center}
\S.\,112.\\[4pt]
\gesperrt{The Principal Resolvent of the Fifth Degree.}
\end{center}

If one makes use of the functions $u, v$ simultaneously, one can derive a family of resolvents which have the form of the principal equation of the fifth degree, and which can be brought directly into agreement with an arbitrarily given principal equation of the fifth degree\footnote{\gesperrt{Kiepert}, G\"ottinger Nachrichten 1878. Crelle's Journal, Bd.\,87. \gesperrt{Klein}, Ikosa\"eder, S.\,106.}.

Since for the icosahedron there are no invariants of degree 8, 14, 16, 22 or 28, the functions $S(\tilde\omega)$, $S(t\tilde\omega)$, $S(\tilde\omega^2)$, $S(t\tilde\omega^2)$, $S(t^2\tilde\omega^2)$ must vanish, and if we therefore set
\weq{1}{Y = \alpha\tilde\omega + \beta t\tilde\omega}
then the functions
\[ S(Y),\quad S(Y^2) \]
will also vanish for arbitrary values of the parameters $\alpha, \beta$, and from this there results an identical equation of the form
%CONT>>
